\section{Discussion}
\label{sec:discussion}
% EDIT (2026-10-29): new section built from the old 5.3 (findings), 6.1--6.2 (gaming discussion,
% detection), 7.4 (what the models changed), 8.1--8.2 (interpretation), 9 (pitfalls, defects) and 10
% (threats to validity), following the reviewer's requested structure. No new number.
This section draws out what the five experiments mean: which prompt components matter, why benchmark
gaming is a benchmark problem rather than a model problem, what the models change and what they miss,
what profiling data is worth, and where the evaluation itself is limited.

\subsection{Which Prompt Components Matter}
\label{sec:analysis}
% EDIT (2026-10-29): restated as a list, as asked in review; every item is read off Tables 1 and 3.
The prompt components are not equally important, and E1 ranks them:
\begin{itemize}\itemsep2pt
  \item \textbf{The correctness argument is the instruction whose removal costs the most.} Removing
    it alone cuts first-attempt success to $8$ of $12$ runs, against $10$--$11$ for every other single
    removal (the \texttt{minimal} prompt, which drops it together with every other instruction,
    still reaches $10$, so the cost comes from keeping the long rule list without the argument that
    anchors it). Removing it is also what lets Gemini game \texttt{shmembench}, and it is the variant
    that gives Gemini its slowest valid result on four of six benchmarks. The requirement costs nothing
    in peak speedup: the best \texttt{bitonic-sort} results of both models are reached with it in
    place.
  \item \textbf{For GPT-5 the instructions matter as a block more than one by one.} No single removal
    hurts it consistently, but stripping the block to one sentence (\texttt{minimal}) gives its slowest
    valid result on four of five benchmarks, by $31$--$34\%$ on two of them.
  \item \textbf{The step-by-step instruction is the safest with respect to slow results and the least
    safe with respect to failures.} \texttt{cot} never gives either model its slowest valid result and
    reaches the best observed speedup for both, but Gemini produced no valid candidate under it on two
    of six benchmarks.
  \item \textbf{The prohibitions and the explanation lists act differently on the two models.}
    Removing the prohibitions produces GPT-5's best \texttt{bitonic-sort} result and one of Gemini's
    worst; removing the explanation lists is Gemini's best variant on the same benchmark and has no
    consistent effect elsewhere. Neither can be recommended or dropped without naming the model.
  \item \textbf{The hardware guidance shows no consistent effect on speed}; its removal produced the
    grid's only timeout, for GPT-5 on \texttt{accuracy}, which Gemini completed without incident.
  \item \textbf{Reliability, not peak speedup, is where models and variants differ most.} GPT-5
    succeeds at the first attempt in $32$ of $36$ runs and Gemini~2.5~Pro in $26$; the best speedup
    either reaches is $1.14\times$.
\end{itemize}

\subsection{Benchmark Gaming and Legitimate Work Removal}
\label{sec:gaming-disc}
% EDIT (2026-10-29): adds the point, asked for in review, that in a generic setting removing
% computation that correct output does not need is good; built from the old 6.2 and the chemv
% paragraph of 8.2.
In a generic setting, the edit that E2 calls gaming is good. A program that produces the correct output
with less computation is a faster program; compilers remove dead code for exactly this reason, and for
an application---whose users care about its answer, not about how much work produced it---removing
work the answer does not need is the goal. \texttt{chemv} (E5) is this case: the model removed loop
structure and barriers that correct output does not require, kept every multiply-add, and the result is
numerically identical to the original.

A benchmark is different because its purpose is the work itself. \texttt{atomicSystemWide} exists to
measure how fast the hardware performs system-scope atomics; a version that performs one atomic
instead of \texttt{loop\_num} of them prints the right total and measures nothing. The same program
transformation is therefore an optimization in one context and a meaningless result in the other,
which is why the semantic-preservation block is scoped to \emph{benchmark scaffolding}---timing,
repeat and convergence loops, verification code, fixed inputs, printed output---and leaves genuine
redundancy inside an application open to elimination. Two further points follow from E1 and E2.
First, a requirement to argue correctness suppresses gaming in one model (Gemini) but not the other
(GPT-5, which gamed \texttt{shmembench} under every variant), so it cannot serve as the defence; the
explicit constraint worked for both. Second, neither output comparison nor the benchmark's own check
can detect this class of edit, by construction, because the exploit preserves the printed result. Our
detection rule is an implausibility threshold---any speedup above $5\times$---followed by reading the
diff. It is cheap enough to run on every candidate, and it also exposed a $31\times$ \texttt{amgmk}
candidate that had deleted the benchmark's report sections. Its limit is equally clear: a gaming
transformation yielding, say, $1.5\times$ would not be flagged, and we cannot rule out such cases in the
reported numbers (Section~\ref{sec:limits}).

\subsection{What the Models Change, and What They Miss}
\label{sec:whatmodels}
% EDIT (2026-10-29): the optimization details moved here from the old 7.4, as asked in review.
The applications show the same pattern from both sides. On Cholla, both GPT-5 and the hand reference
replace $18$ double-precision divisions by $6$ reciprocal multiplications; GPT-5 additionally removes
integer index arithmetic that was used only for a bounds check, which reduces register pressure from
$109$ to $89$ registers per thread without spilling.
These are local, exact transformations: each is confined to a few lines, each is provably equivalent,
and each is the kind of change the models found reliably across benchmarks as well---reciprocal
reuse, redundant index arithmetic removed, read-only loads. On AutoDock-GPU the profitable change is
of a different kind. The hand reference fuses seven sequential block reductions, each with three
barriers and a shared-memory atomic, into one pass, worth $1.532\times$; every candidate that built
left all seven reductions in place and gained at most $3.7\%$. Restructuring across a whole kernel is
what the models miss, and giving them more freedom did not help: the structural prompt never
outperformed the baseline on either application.

The three real-application failures share one cause: the model cannot see how its file is compiled,
whether it is a standalone translation unit or textually included, which symbols are already in scope,
and which declarations other files depend on. This never arises for a self-contained benchmark and
is, in our experience, the dominant failure mode on real code. It also means some optimizations a human
would make---changing a header, a launch site or a build flag---are out of reach by construction. Two
directions follow. The prompt should carry the file's \emph{compilation context}, which is cheap to
extract and addresses this failure mode directly; and the structural freedom that helps on standalone
kernels needs a real-application contract that forbids changing externally visible signatures while
still allowing file-internal restructuring.

\subsection{Profiling Data, Raw and Interpreted}
\label{sec:profdisc}
E4's pattern is consistent with the workloads rather than with the models. A HeCBench kernel is a few
hundred lines with one obvious hot loop, and the profile tells a model little it cannot infer from the
source: it changes which edit the model makes, sometimes by more than $10\%$ either way, but it does not
make the edit better on average. Cholla's HLLC solver is $307$ lines inside a 21k-line CUDA application
whose cost is spread over several kernels; there the profile identifies both the kernel worth editing
and the reason it is slow (latency-bound at $19.7\%$ DRAM and $40.4\%$ SM throughput, occupancy $23\%$
limited by $109$ registers per thread). We therefore expect profile data to pay off in proportion to
how hard the bottleneck is to locate from source alone, which is the regime real applications occupy.
With one application showing the gain, this is a hypothesis rather than a measured law.

E5 shows what happens when the data is interpreted for the model. A checklist of remedies is an anchor
as much as a hint: when the listed transformation is the best available one, anchoring is worth a large
factor (\texttt{bincount}, $1.00\times \to 1.38\times$ for both models); when it is not, the model
applies it anyway (\texttt{concat}, where GPT-5 took ``several elements per thread'' to one thread per
$128$-element line). The mean over twelve pairs is $0.97\times$, which is why a mean speedup, the
statistic the literature reports, hides the effect; the result is the dispersion. Two practical
consequences follow. Prescriptive guidance should be reported with its variance, not only its mean.
And the guidance should be \emph{evidence rather than remedy}: state the measured condition (``$23\%$
occupancy, limited by $109$ registers per thread'') and let the model decide what follows, instead of
naming the transformation. We did not get to test that variant, and it is the first thing we would run
next.

\subsection{Benchmarks Versus Applications}
\label{sec:bva}
The benchmark results understate how much prompt design matters, because single-kernel benchmarks hide
two of its effects. They cannot expose compilation-context failures, since each benchmark is a
self-contained translation unit; and they make profiling data look worthless, since the hot loop of a
$200$-line benchmark is obvious from the source. Both effects reappear on the production applications,
which is the strongest argument we can offer against drawing conclusions about such systems from
benchmark suites alone.

\subsection{Evaluation Pitfalls and Tooling Defects}
\label{sec:pitfalls}
\label{sec:defects}
Running the pipeline on production code exposed problems that benchmark runs had hidden, some of them
in the evaluation itself. Table~\ref{tab:pitfalls} collects them with the practice we adopted in
response.

% NEW: summary table. Every row restates a finding made elsewhere in the paper; no new measurement.
\begin{table}[t]
\centering
\footnotesize
\caption{Evaluation pitfalls observed in this study, where each is discussed, and the practice adopted.}
\label{tab:pitfalls}
\begin{tabular}{@{}p{0.27\textwidth}p{0.10\textwidth}p{0.52\textwidth}@{}}
\toprule
\textbf{Pitfall} & \textbf{Section} & \textbf{Practice adopted} \\
\midrule
Edits that delete the measured work while preserving output and the built-in check
  & \ref{sec:intent}, \ref{sec:gaming-disc} & Semantic-preservation constraint; any speedup above $5\times$ is a correctness alarm and its diff is read \\
Driver update on a shared node: wrong output and inflated times, with no error
  & \ref{sec:limits} & Re-time the unmodified program at the start of every session and compare with earlier measurements of the same binary \\
Built-in check that cannot pass for the unmodified program, or prints no verdict
  & \ref{sec:correctness}, \ref{sec:guidvar} & Run the unmodified program through the check first; exclude benchmarks whose check cannot pass; explicit gates for applications \\
Drift between sessions on shared nodes (up to $49\%$ on the same binary)
  & \ref{sec:timing} & Build once, interleave arms, wait for an idle GPU, log foreign GPU processes per repetition, require repetitions to agree \\
Application output with no built-in check, or not reproducible run to run
  & \ref{sec:correctness} & One gate per application, bit-exact or a statistical band; validate each gate by fault injection \\
Wall clock dominated by setup (AutoDock-GPU CUDA setup: $68$\,s)
  & \ref{sec:timing} & Use each application's own phase timer \\
Prompt contents differing from what was believed sent
  & \ref{sec:defects} & Defects fixed; earlier results reported as measured and not regenerated \\
Single samples and noisy pairs
  & \ref{sec:timing}, \ref{sec:limits} & Score a pair a tie when its difference is inside its own spread; treat sub-$5\%$ single-cell effects as unresolved \\
Compilation context invisible to a single-file editor
  & \ref{sec:whatmodels} & Reported as the dominant real-code failure mode; carrying that context in the prompt is future work \\
\bottomrule
\end{tabular}
\end{table}

\paragraph{Prompt-fidelity defects exposed by real code.}
Running production applications through the loop surfaced a class of problem that benchmark runs had
hidden for months: the prompt the model actually received was not the prompt we believed we were
sending. Table~\ref{tab:defects} lists the defects we found and fixed. They fall into two groups.
\emph{Summarization} defects corrupt the profiling evidence---a shared-memory footprint of
$7.43$\,Kbyte presented to the model as ``7 bytes'', SI \texttt{Kbyte} read as $1024$ bytes, and, most
seriously, a CSV reader that parsed every section of an Nsight Compute source page against the first
section's header and silently dropped rows of any other width. On one benchmark that left $36$ of $76$
stall rows in the prompt; re-running $56$ profiled inputs after the fix changed four of them, all by
\emph{gaining} rows. \emph{Loop} defects corrupt the repair channel: a build timeout reported to the
model as a fault in its own code, and an error filter that kept only lines matching
\texttt{error}$|$\texttt{fatal}$|$\texttt{undefined reference}, so that a linker's explanatory line
(\texttt{multiple definition of \ldots first defined here}) was discarded and the model saw only
\texttt{ld returned 1 exit status}.

One defect deserves separate mention because it is a design decision rather than an oversight. The
pipeline strips all comments from the source before prompting, to save tokens, and then replaces the
original file with the model's full-file answer---so every optimized file lost its license header and
all of its documentation. We initially, and wrongly, attributed the missing LGPL header of AutoDock-GPU
to the models; the prompts contained no occurrence of ``Copyright''. We now re-attach the original
leading comment block when it is a license and the candidate lacks it, but the general question stands:
comment stripping also discards institutional knowledge such as \texttt{// DO NOT UNDEFINE} directives
that exist precisely to stop the edit an optimizer is about to make.

\begin{table}[t]
\centering
\footnotesize
\setlength{\tabcolsep}{3pt}
\caption{Prompt-fidelity defects found while running production applications, and what each one
  corrupted. All were present during the benchmark study.}
\label{tab:defects}
\begin{tabular}{@{}p{0.30\columnwidth}p{0.63\columnwidth}@{}}
\toprule
\textbf{Defect} & \textbf{Effect} \\
\midrule
Shared memory read without unit & $7.43$\,Kbyte/block shown as ``7 bytes''; static and dynamic sizes
  summed across units \\
\texttt{Kbyte} treated as $1024$\,B & $2.4\%$ error in every stated size \\
Per-section CSV headers ignored & $25$ of $26$ sections dropped on one input ($36/76$ stall rows
  reached the model) \\
Stall reasons from a Python \texttt{set} & Prompt text not reproducible across processes \\
Kernel names matched verbatim & Namespaced and templated kernels fell back to the capped launch count \\
Comments stripped before prompting & Optimized files lost their license header and all documentation \\
Echoed line numbers not stripped & Uncompilable output when a model echoes the numbered source \\
Fixed $300$\,s build timeout & Real-application rebuilds reported to the model as its own fault \\
Linker cause lines filtered out & Repair prompt contained no actionable error \\
\bottomrule
\end{tabular}
\end{table}

\subsection{Limitations and Threats to Validity}
\label{sec:limits}
% EDIT (2026-10-29): condensed from the old Section 10; the measurement protocol and the gate
% definitions are stated in Section 2 and not repeated.
\paragraph{Measurement.}
Even under the protocol of Section~\ref{sec:timing}, residual interference cannot be excluded on a
shared machine: in one Cholla pass three repetitions drifted by $15\%$ while our detector reported no
foreign GPU process. Measurement noise alone puts a floor of roughly $2\%$ under any effect we report on
the applications, and far higher on the noisiest benchmarks, where per-repetition spread exceeds
$30\%$. The benchmark half of E4 was first timed after a driver update on the shared H100 node, which
made binaries built with CUDA 12.x print wrong output and run $20$--$300\times$ slower without failing;
unmodified benchmarks then printed FAIL from their own checks and two exceeded our $600$\,s per-run cap.
We detected this only by comparing against earlier measurements of the same binaries, discarded that
timing, and re-timed the same candidates rebuilt with CUDA 13.3.1; Section~\ref{sec:profabl} reports
only the re-timing. The candidates themselves were generated on the faulty node, so a candidate
rejected by the generation-time check could have been rejected wrongly; the one case where this decides
a result (\texttt{background-subtract}, GPT-5) we re-checked with the working toolchain, and the
rejection stands.

\paragraph{Tooling defects.}
The defects of Table~\ref{tab:defects} were all present during E1: profiling summaries reached the
model with wrong units, truncated stall data and understated launch counts, and the repair channel
discarded the diagnostics a model needed. They are fixed, but every result predating a fix was produced
with the affected prompt. We report E1 as measured rather than re-running the grid, and we do not claim
that it would reproduce unchanged on the corrected summarizers.

\paragraph{Sampling of model outputs.}
Except where stated, each (program, model, prompt) cell is a single LLM sample, so a cell conflates the
prompt's effect with sampling variance. Only one setting is replicated: a second, independently
generated sample of the baseline prompt on both applications, which keeps Cholla on the same side of
the profile-data comparison and reproduces AutoDock-GPU's tie. The prescriptive-guidance comparison
(E5) has no valid replicate: a second generation round was timed on the H100 node during the fault
described above and was stopped after six pairs. Combined with the measurement floor, an effect below
roughly $5\%$ in a single cell should be treated as unresolved rather than small. E4 spans $24$ pairs
and scores $23$; its positive result rests on a single application, Cholla, with two models, and the
nineteen benchmark pairs show that profile data moves results in both directions there, not that the
gain on Cholla generalizes to other large applications.

\paragraph{Correctness is empirical.}
The gates of Section~\ref{sec:correctness} were validated by fault injection, but a subtle accuracy
loss that stays inside the AutoDock-GPU band would pass, and the gates check printed outputs rather
than all internal state. For the three benchmarks that print no verdict (\texttt{amgmk},
\texttt{atomicSystemWide}, \texttt{bincount}), correctness rests on the pipeline's generation-time
check and on our reading of the diffs, which is weaker than the application gates. Gaming is detected
heuristically (Section~\ref{sec:gaming-disc}): the $5\times$ alarm catches every case we observed,
all of them above $30\times$, but would miss a small one.

\paragraph{Scope.}
We optimize the execution time of one GPU kernel file under a correctness gate; energy, memory
footprint, multi-GPU scaling, numerical stability over long runs and the maintainability of the
generated code are not studied. Results cover NVIDIA GPUs only, one architecture per study, so an
effect that depends on register file size, SM count or cache hierarchy would be invisible to us, and
the E1 profiles were collected on an A100 and used on an H100. We use two proprietary models whose
weights and versions we do not control; the same prompt may behave differently after a model update,
and the Gemini version changed between E1 and the later experiments.
